\chapter{The simplified controller}
\label{ch:kursen}

{\large\itshape What a teaching controller implements, and what it deliberately leaves out.\par}

\begin{chapterbox}{In this chapter}
\begin{itemize}
  \item The blocks a simplified CAN controller is made of.
  \item Which parts of \kapref{ch:buss}--\kapref{ch:fel} are included, and which are not.
  \item Why polling software cannot talk to the controller directly.
  \item The register layer that makes it reachable, and the three limitations that reach all the
    way up into the code.
\end{itemize}
\end{chapterbox}

The six preceding chapters describe CAN as the standard defines it. This chapter describes the
controller that is actually built in teaching: what it does of all that, and what it leaves out.

The chapter faces both ways. If you build the controller, the list below is the requirements
specification. If you write the driver, it is instead the list of what you must not count on.

\section{The blocks}

\bookfigure{controller_blocks}{The simplified controller's blocks, and the layer that makes it
reachable from a processor.}{fig:blocks}

\begin{booktable}{@{}lL@{}}
\thead{Block} & \thead{Responsibility} \\ \midrule
Bit timer & Counts ticks per bit, pulses at the sample point and at the end of the bit period
  (\kapref{ch:tajming}). \\
CRC-15 & A bit-serial engine that both computes and checks (\kapref{ch:stuffing}). \\
Transmit shift register & Serialises MSB first and inserts stuff bits. \\
Receive shift register & Deserialises and removes stuff bits; reports stuff errors. \\
Controller & The state machine that sequences the frame's fields and handles arbitration. \\
Synchroniser & Two flip-flops on every asynchronous input, so that the bus and reset can be read
  safely. \\
\end{booktable}

On top of those sit two more blocks, which have nothing to do with CAN but with making the
controller \emph{reachable}:

\begin{booktable}{@{}lL@{}}
\thead{Block} & \thead{Responsibility} \\ \midrule
Register bank & Turns the controller's single-cycle pulses into sticky, pollable registers. \\
Serial bridge & Makes the registers readable and writable from outside, over a serial bus. \\
\end{booktable}

\section{What is included}

The controller implements the frame's sequence in full: SOF, the arbitration field with identifier
and RTR, the control field with IDE, r0 and DLC, the data field, CRC-15 with delimiter, ACK slot
with delimiter, and a seven-bit EOF.

It also implements:

\begin{itemize}
  \item \textbf{Bit stuffing in both directions}, with the stuffing state kept alive across field
    boundaries.
  \item \textbf{CRC-15}, with the same engine on the transmit and receive sides.
  \item \textbf{Arbitration}, bit by bit across the arbitration field.
  \item \textbf{Open drain}: the node drives only when it should, and the bus is dominant if anyone
    pulls it down.
  \item \textbf{Stuff error detection} and \textbf{CRC checking} on the receive side.
\end{itemize}

That is enough for two such nodes to exchange real CAN frames on a real bus, and enough for a
commercial bus analyser to decode them.

\section{What is not included}

\begin{booktable}{@{}lL@{}}
\thead{Missing} & \thead{Consequence} \\ \midrule
Error frames & An error does not abort the transfer for the other nodes. \\
Error counters & No error passive, no bus off (\kapref{ch:fel}). \\
Automatic retransmission & A frame that went wrong is not resent by the hardware. \\
ACK readback & The transmitter never learns whether anyone acknowledged. \\
Reception while transmitting & A node that loses arbitration does not receive the frame that won. \\
Receive queue & The controller holds exactly one received frame. \\
Filtering & Every received frame is passed on; filtering is done in software. \\
Extended identifiers & 11 bits only. \\
Bit error detection outside arbitration & The transmitter reads the bus back only during the arbitration field. \\
\end{booktable}

Each of them is a deliberate choice. Together they are the difference between a design that can be
built and verified during a course, and one that cannot.

\section{Why a register layer is needed}

The controller is built for a caller in the same clock domain. It signals with pulses that are
\textbf{one clock cycle} long --- 20~ns at 50~MHz.

A processor polling over a serial bus sees the system a few tens of microseconds at a time at best.
A 20~ns pulse is therefore gone thousands of cycles before the next poll comes.

\textbf{The register bank solves that by turning pulses into levels.} A pulse sets a bit, and the
bit stays set until the software explicitly clears it:

\begin{booktable}{@{}lll@{}}
\thead{Status bit} & \thead{Set by} & \thead{Cleared by} \\ \midrule
TX done & a completed transmission & an accepted transmit trigger \\
RX valid & a received frame & an acknowledgement from the software \\
Error & a rising edge on the error level & an explicit clear \\
\end{booktable}

That pattern --- \emph{sticky bit, explicit clear} --- is the whole point of the layer, and it is
also the most common source of bugs in the software above it. Forget the clear and the bit stays set
forever, and the driver reads the same frame in an endless loop.

\begin{aside}
\note The register bank does not \emph{hide} the controller's limitations, it \emph{passes them on}.
That is a design choice worth noticing: a layer that buffered received frames would have looked
helpful and made the system's real behaviour impossible to reason about.
\end{aside}

\section{The three limitations that reach all the way up}

Three of the simplifications above are not merely internal. They show in the software, and whoever
writes the driver has to handle them rather than be surprised by them.

\subsection{A lost arbitration looks like a successful transmission}

The controller signals ``the transmission is over'' in both cases. The bit means \textbf{the port is
free again}, not \emph{the frame got through}.

That it behaves this way is deliberate, and it is in fact the fix for a real bug: if it signalled
only on a successful transmission, the bit would be stuck after every lost arbitration, and the
node could not send again until it was reset by hand. On a two-node bus, that is normal
operation.

The consequence for the software is a polling loop in two steps:

\begin{console}
wait for "port is free"   -> the transmit attempt is over
read the error flag       -> did it go well or not?
\end{console}

Without the second step there is no way to know whether the frame should be resent.

\subsection{One bit for four causes}

Lost arbitration, stuff error, failed CRC and received remote frame all give the same single error
flag. The software can see \emph{that} something went wrong, never \emph{what}.

Retry logic that treats a lost arbitration differently from a broken CRC therefore cannot be written
against this interface. The right handling is to know about the limitation and treat all four
alike --- not to guess the cause from the context.

\subsection{No buffering on the receive side}

The controller holds one frame. If the next one arrives before the previous one has been
acknowledged, it is overwritten, and the newest wins.

That is documented and deliberate. A real controller has a receive queue; this one does not, and
the software has to poll often enough. At some traffic intensity that is not enough however often it
polls, and frames are then dropped --- which is the hardware's limitation and not the software's
bug.

\section{What a real controller would add}

For completeness, and as the answer to the question ``how far is there left to go?'':

\begin{itemize}
  \item Error frames, error counters and bus off (\kapref{ch:fel}).
  \item Automatic retransmission after errors.
  \item Receive filters in hardware, so that uninteresting frames never reach the software.
  \item A receive queue with several slots.
  \item Several transmit buffers with priorities among them.
  \item Extended identifiers, and with them CAN 2.0B.
  \item Readback of the whole frame, and with it real bit error detection.
\end{itemize}

The list is long, but notice what is \emph{not} on it: framing, stuffing, CRC, arbitration and bit
timing. The hard part of CAN is already done.

\section*{Exercises}

\exercise{What is missing}{Reasoning}
\label{ex:7:1}
For each of the following scenarios: what does a complete CAN controller do, and what does the
simplified one do?
\begin{enumerate}[a)]
  \item A frame loses arbitration.
  \item A received frame has the wrong checksum.
  \item No node acknowledges a transmitted frame.
  \item Two frames are received in quick succession.
\end{enumerate}

\exercise{Sticky bits}{Reasoning}
\label{ex:7:2}
The controller pulses for one clock cycle; the software polls every fortieth microsecond.
\begin{enumerate}[a)]
  \item How many 50~MHz clock cycles pass between two polls?
  \item What would the software see if the status bit mirrored the pulse directly?
  \item Why must the bit be cleared by a write instead of automatically on read?
\end{enumerate}

\exercise{Two steps}{Reasoning}
\label{ex:7:3}
A driver waits for the transmission to complete and then returns ``it went well''.
\begin{enumerate}[a)]
  \item What mistake does it make?
  \item What must it do instead?
  \item Which four causes can lie behind a set error flag?
\end{enumerate}

\exercise{The boundary}{Reasoning}
\label{ex:7:4}
The register bank passes the controller's limitations on instead of hiding them.
\begin{enumerate}[a)]
  \item Give an example of how it \emph{could} hide one of them.
  \item What would that have cost whoever debugs the system?
  \item What principle does this say something about, more generally, when designing an interface
    between hardware and software?
\end{enumerate}
